فرض کنید تصویر اصلی $f(x,y)$ با اندازه $M\times N$ داریم و هر دو روش، آن را به یک اندازهٔ یکسان $P\times Q$ صفرپَد می‌کنند (پس تعداد صفرها برابر است). تفاوت اصلی این دو روش این است که مکانِ قرارگیریِ $f$ داخل آرایهٔ پَدشده عوض می‌شود:

\begin{itemize}
\item در پَد نامتقارن (گوشه‌ای)، تصویر در گوشهٔ بالا-چپ قرار می‌گیرد؛ یعنی عملاً شیفتی ندارد.
\item در پَد متقارن (مرکزی)، تصویر به اندازهٔ $(x_0,y_0)$ داخل قابِ بزرگ‌تر شیفت داده می‌شود تا وسط قرار بگیرد.
\end{itemize}

اگر $p(x,y)$ نسخهٔ پَدشده باشد، وقتی فقط مکانِ تصویر داخل قاب عوض شود داریم:
$$
p_{\text{sym}}(x,y)=p_{\text{asym}}(x-x_0,\;y-y_0).
$$
خاصیت شیفت در \lr{DFT} می‌گوید:
$$
P_{\text{sym}}(u,v)=P_{\text{asym}}(u,v)\;e^{-j2\pi\left(\frac{ux_0}{P}+\frac{vy_0}{Q}\right)}.
$$
پس این دو تبدیلِ فوریه فقط در یک فاکتورِ فاز (وابسته به فرکانس) فرق دارند و اندازهٔ طیفشان یکی است:
$$
|P_{\text{sym}}(u,v)|=|P_{\text{asym}}(u,v)|.
$$

حال در فیلترکردنِ حوزهٔ فرکانس، $G(u,v)=H(u,v)P(u,v)$.
بنابراین:
$$
G_{\text{sym}}(u,v)=G_{\text{asym}}(u,v)\;e^{-j2\pi\left(\frac{ux_0}{P}+\frac{vy_0}{Q}\right)}
\;\;\Longrightarrow\;\;
g_{\text{sym}}(x,y)=g_{\text{asym}}(x-x_0,\;y-y_0).
$$

\subsection*{نتیجه}
این دو نوع پَد کردن (با اندازهٔ پَد یکسان) نتیجهٔ «بنیاداً متفاوتی» تولید نمی‌کنند؛ خروجی‌ها معادل‌اند ولی نسبت به هم فقط یک جابه‌جایی مکانی (و در حوزهٔ فرکانس یک فازِ خطی) دارند. اگر بعد از \lr{IDFT} ناحیهٔ درست را با همان مرجع مکانی \lr{crop} کنید، به یک نتیجهٔ یکسان می‌رسید؛ در غیر این صورت خروجی ظاهراً شیفت‌خورده دیده می‌شود.